\subsection{How Sequences Are Specified}
A sequence is an ordered list of numbers, usually written as $(a_n)_{n\ge 1}$.
Common ways to define a sequence are:
\begin{itemize}[leftmargin=*]
  \item explicit formula: $a_n$ written directly in terms of $n$
  \item recursive formula: each term written using previous terms
  \item contextual rule: terms generated from a worded process
\end{itemize}

\textit{Examples.}
\[
a_n = 3n-1,\quad b_1=2,\ b_{n+1}=b_n+5,\quad c_n = (-1)^n.
\]

\subsection{Arithmetic Sequences (AP)}
An arithmetic sequence has constant difference $d$:
\[
a_{n+1}-a_n=d.
\]
If the first term is $a$, then
\[
a_n = a + (n-1)d.
\]

\begin{problem}[Recognizing an AP]
Given $a_1=7$ and $a_{n+1}=a_n-3$, find $a_{10}$.
\end{problem}
\begin{solution}
Here $d=-3$, so
\[
a_{10}=7+(10-1)(-3)=7-27=-20.
\]
\end{solution}

\subsection{Arithmetic Series}
The sum of the first $n$ terms of an AP is:
\[
S_n = \frac{n}{2}\bigl(2a+(n-1)d\bigr)=\frac{n}{2}(a+l),
\]
where $l$ is the $n$th term.

\begin{problem}[AP Sum]
Find $\sum_{k=1}^{20}(5+2(k-1))$.
\end{problem}
\begin{solution}
This is an AP with $a=5$, $d=2$, and $n=20$. The 20th term is
\[
l=5+19\cdot 2=43.
\]
So
\[
S_{20}=\frac{20}{2}(5+43)=10\cdot 48=480.
\]
\end{solution}

\subsection{Geometric Sequences (GP)}
A geometric sequence has constant ratio $r$:
\[
\frac{a_{n+1}}{a_n}=r \quad (a_n\ne 0).
\]
If the first term is $a$, then
\[
a_n = ar^{n-1}.
\]

\begin{problem}[Recognizing a GP]
If $a_1=81$ and $r=\frac13$, find $a_5$.
\end{problem}
\begin{solution}
\[
a_5 = 81\left(\frac13\right)^4 = 81\cdot\frac{1}{81}=1.
\]
\end{solution}

\subsection{Geometric Series}
For $r\ne 1$, the sum of the first $n$ terms is
\[
S_n = a\frac{1-r^n}{1-r} = a\frac{r^n-1}{r-1}.
\]
Use the form that keeps signs clean for your $r$ value.

\begin{problem}[Finite GP Sum]
Find $\sum_{k=0}^{6}3\cdot 2^k$.
\end{problem}
\begin{solution}
This is $a=3$, $r=2$, $n=7$ terms (from $k=0$ to $6$):
\[
S_7 = 3\frac{2^7-1}{2-1}=3(128-1)=381.
\]
\end{solution}

\subsection{Limiting Sum of a Geometric Series}
If $|r|<1$, then $r^n\to 0$ as $n\to\infty$, and the infinite sum exists:
\[
S_\infty = \frac{a}{1-r}.
\]
If $|r|\ge 1$, the geometric series does not converge to a finite limiting sum.

\begin{problem}[Infinite GP]
Evaluate $2+\frac{2}{3}+\frac{2}{9}+\frac{2}{27}+\cdots$.
\end{problem}
\begin{solution}
Here $a=2$, $r=\frac13$, and $|r|<1$, so
\[
S_\infty=\frac{2}{1-\frac13}=\frac{2}{\frac23}=3.
\]
\end{solution}

\subsection{Recurring Decimals and Geometric Series}
Recurring decimals can be written as geometric series.
For example,
\[
0.\overline{7}=0.7+0.07+0.007+\cdots
\]
is geometric with $a=0.7$ and $r=0.1$. So
\[
0.\overline{7}=\frac{0.7}{1-0.1}=\frac{0.7}{0.9}=\frac79.
\]

\subsection{Sigma Notation Quick Rules}
\[
\sum_{k=1}^{n}(u_k+v_k)=\sum_{k=1}^{n}u_k+\sum_{k=1}^{n}v_k,
\quad
\sum_{k=1}^{n}cu_k = c\sum_{k=1}^{n}u_k.
\]
Index shifts must preserve both the expression and bounds consistently.

\begin{takeaways}
\item Always identify whether the question asks for a term $a_n$, a finite sum $S_n$, or a limit.
\item AP questions depend on constant difference; GP questions depend on constant ratio.
\item For infinite geometric sums, check $|r|<1$ before using $S_\infty=\frac{a}{1-r}$.
\item Track indexing carefully in sigma notation to avoid off-by-one errors.
\item Chebyshev polynomial questions usually depend on either the recurrence relation or the trigonometric identities.
\end{takeaways}
